\section{Conclusion}

The Attachment Regulation and Control (ARC) framework provides a parsimonious, empirically testable bridge between adult attachment dispositions and moment-to-moment relational regulation. It does not claim to invent feedback control in attachment theory. Its proposed contribution is to integrate established ideas into a temporal measurement architecture for adult romantic relationships:
\[
\hat{\mathbf r}_t,\mathbf r^\ast_t
\rightarrow
\mathbf e_t
\rightarrow
\mathbf a_{t+1}
\rightarrow
\mathbf p_{t+1}
\rightarrow
\hat{\mathbf r}_{t+1}
\rightarrow
\mathrm{recovery}.
\]

The framework makes two levels of claim independently falsifiable. First, desired and perceived relational state must contain information about subsequent regulation beyond context and behavioral inertia. Second, the parsimonious difference representation \(e_t=r^\ast_t-\hat r_t\) must compete successfully with perceived-state-only, unconstrained component, and asymmetric discrepancy alternatives. Failure of the second claim revises the mathematical representation without automatically rescuing or rejecting the broader feedback account.

Attachment anxiety and avoidance are then tested as moderators of channel-specific regulation rather than as fixed controller labels. Positive secure-base representation is deliberately kept outside the required core until H3 demonstrates reproducible incremental recovery value beyond low anxiety and avoidance and, in stronger tests, beyond proximal partner support. Mentalized affectivity and emotion-regulation flexibility remain optional moderators. Adult disorganization remains an external criterion.

ARC earns explanatory value only if these dynamic measurements improve reproducible, held-out prediction beyond simpler alternatives. Its first decisive test is therefore empirical rather than computational: whether repeatedly measured relational states and targets predict subsequent regulatory action in human relationships, and whether the proposed discrepancy representation is actually warranted by the data.
